% GENERATED FILE. Edit canonical lessons, then run npm run generate:books.
% canonical-source-hash: fnv1a64:416e0f7527082756
% canonical-lessons: TA-C274-icai, TA-C274-patal, TA-C274-katai, TA-C274-katavul, TA-C274-vilanku

\chapter{Music, Song, Story, God, Animal}
\label{ch:ta-music-song-story-god-animal}
\hlchaptermodality{274}

\begin{chapteropening}
\textbf{By the end of this chapter:} \emph{I can say music, a song, a story, God and an animal.}

Complete the last lesson of chapter 274: I can say music, a song, a story, God and an animal.
\end{chapteropening}

\section[\texorpdfstring{\ta{இசை} (icai)}{இசை (icai)}]{\ta{இசை} (icai) — music}

\label{lesson:TA-C274-icai}



\begin{quote}
Before the new one: say the Tamil for twins, then the Tamil for a brother-in-law.
\end{quote}

\subsection*{You'll want to know: \ta{இசை}}
\textbf{\ta{இசை}} — \emph{icai} — \textquotedblleft{}music\textquotedblright{}.

\begin{grammarlens}[title={What you've built}]
One more word for this chapter.
\end{grammarlens}

\subsection*{Your turn}
\begin{itemize}
\raggedright
  \item \emph{Say it:} \emph{icai}
  \item \emph{Say it:} \emph{icai}, once more
  \item \emph{Say it:} say \emph{icai}
  \item \emph{Recall:} say the Tamil for twins, then the Tamil for a brother-in-law, then say \emph{icai} again
\end{itemize}

\begin{tcolorbox}[breakable,colback=teal!4,colframe=teal!35!black,title={Before you move on}]
What does \ta{இசை} mean? (\textquotedblleft{}Music\textquotedblright{}.) Say it once more.
\end{tcolorbox}
\section[\texorpdfstring{\ta{பாடல்} (pāṭal)}{பாடல் (pāṭal)}]{\ta{பாடல்} (pāṭal) — a song}

\label{lesson:TA-C274-patal}



\begin{quote}
Before the new one: say the Tamil for a brother-in-law, then the Tamil for music.
\end{quote}

\subsection*{You'll want to know: \ta{பாடல்}}
\textbf{\ta{பாடல்}} — \emph{pāṭal} — \textquotedblleft{}a song\textquotedblright{}.

\begin{grammarlens}[title={What you've built}]
One more word for this chapter.
\end{grammarlens}

\subsection*{Your turn}
\begin{itemize}
\raggedright
  \item \emph{Say it:} \emph{pāṭal}
  \item \emph{Say it:} \emph{pāṭal}, once more
  \item \emph{Say it:} say \emph{pāṭal}
  \item \emph{Recall:} say the Tamil for a brother-in-law, then the Tamil for music, then say \emph{pāṭal} again
\end{itemize}

\begin{tcolorbox}[breakable,colback=teal!4,colframe=teal!35!black,title={Before you move on}]
What does \ta{பாடல்} mean? (\textquotedblleft{}A song\textquotedblright{}.) Say it once more.
\end{tcolorbox}
\section[\texorpdfstring{\ta{கதை} (katai)}{கதை (katai)}]{\ta{கதை} (katai) — a story}

\label{lesson:TA-C274-katai}



\begin{quote}
Before the new one: say the Tamil for music, then the Tamil for a song.
\end{quote}

\subsection*{You'll want to know: \ta{கதை}}
\textbf{\ta{கதை}} — \emph{katai} — \textquotedblleft{}a story\textquotedblright{}.

\begin{grammarlens}[title={What you've built}]
One more word for this chapter.
\end{grammarlens}

\subsection*{Your turn}
\begin{itemize}
\raggedright
  \item \emph{Say it:} \emph{katai}
  \item \emph{Say it:} \emph{katai}, once more
  \item \emph{Say it:} say \emph{katai}
  \item \emph{Recall:} say the Tamil for music, then the Tamil for a song, then say \emph{katai} again
\end{itemize}

\begin{tcolorbox}[breakable,colback=teal!4,colframe=teal!35!black,title={Before you move on}]
What does \ta{கதை} mean? (\textquotedblleft{}A story\textquotedblright{}.) Say it once more.
\end{tcolorbox}
\section[\texorpdfstring{\ta{கடவுள்} (kaṭavuḷ)}{கடவுள் (kaṭavuḷ)}]{\ta{கடவுள்} (kaṭavuḷ) — God, a god}

\label{lesson:TA-C274-katavul}



\begin{quote}
Before the new one: say the Tamil for a song, then the Tamil for a story.
\end{quote}

\subsection*{You'll want to know: \ta{கடவுள்}}
\textbf{\ta{கடவுள்}} — \emph{kaṭavuḷ} — \textquotedblleft{}God, a god\textquotedblright{}.

\begin{grammarlens}[title={What you've built}]
One more word for this chapter.
\end{grammarlens}

\subsection*{Your turn}
\begin{itemize}
\raggedright
  \item \emph{Say it:} \emph{kaṭavuḷ}
  \item \emph{Say it:} \emph{kaṭavuḷ}, once more
  \item \emph{Say it:} say \emph{kaṭavuḷ}
  \item \emph{Recall:} say the Tamil for a song, then the Tamil for a story, then say \emph{kaṭavuḷ} again
\end{itemize}

\begin{tcolorbox}[breakable,colback=teal!4,colframe=teal!35!black,title={Before you move on}]
What does \ta{கடவுள்} mean? (\textquotedblleft{}God, a god\textquotedblright{}.) Say it once more.
\end{tcolorbox}
\section[\texorpdfstring{\ta{விலங்கு} (vilaṅku)}{விலங்கு (vilaṅku)}]{\ta{விலங்கு} (vilaṅku) — an animal}

\label{lesson:TA-C274-vilanku}



\begin{quote}
Before the new one: say the Tamil for a story, then the Tamil for God.
\end{quote}

\subsection*{You'll want to know: \ta{விலங்கு}}
\textbf{\ta{விலங்கு}} — \emph{vilaṅku} — \textquotedblleft{}an animal\textquotedblright{}.

\begin{grammarlens}[title={What you've built}]
One more word for this chapter.
\end{grammarlens}

\subsection*{Your turn}
\begin{itemize}
\raggedright
  \item \emph{Say it:} \emph{vilaṅku}
  \item \emph{Say it:} \emph{vilaṅku}, once more
  \item \emph{Say it:} say \emph{vilaṅku}
  \item \emph{Recall:} say the Tamil for a story, then the Tamil for God, then say \emph{vilaṅku} again
\end{itemize}

\begin{tcolorbox}[breakable,colback=teal!4,colframe=teal!35!black,title={Before you move on}]
What does \ta{விலங்கு} mean? (\textquotedblleft{}An animal\textquotedblright{}.) Say it once more.
\end{tcolorbox}
